\section*{Bab \ref{ch:b2:climate-proxies} --- Arsip dan Proksi Iklim}

\begin{solution}{exo:b2:climate-proxies:1}
(a) $\delta = (27.5 \times (-30) + 1350 \times 0)/(1350 + 27.5) = -825/1377.5 = -0.60$ permil.
(b) Samudra lalu akan memuat seluruh air: $-0.60$ permil, \num{0.60} permil lebih ringan daripada
masa kini. (c) $27.5/1377.5 = \qty{2.0}{\percent}$.
\end{solution}

\begin{solution}{exo:b2:climate-proxies:2}
(a) $T = 16.1 - 4.76\,(2.50 - (0 - 0.27)) = 16.1 - 4.76 \times 2.77 = 16.1 - 13.19 =
\qty{2.9}{\celsius}$. (b) $T = 16.1 - 4.76\,(2.50 - 0.73) = 16.1 - 8.43 = \qty{7.7}{\celsius}$:
cangkang yang sama memberikan suhu hampir lima derajat lebih tinggi jika airnya lebih berat, dan karena itulah
komposisi air harus diketahui. (c) Cangkang itu $2.71 - 2.50 = 0.21$ permil lebih berat,
sehingga airnya $4.76 \times 0.21 = \qty{1.0}{\celsius}$ lebih dingin, pada \qty{1.9}{\celsius}.
\end{solution}

\begin{solution}{exo:b2:climate-proxies:3}
(a) $0.4\,\mathrm{e}^{0.9} = \qty{0.98}{mmol/mol}$ dan $0.4\,\mathrm{e}^{1.8} =
\qty{2.42}{mmol/mol}$. (b) Faktornya $\mathrm{e}^{0.9} = 2.46$; per derajat $\mathrm{e}^{0.09}
= 1.094$, pertumbuhan sebesar \qty{9.4}{\percent} per derajat. (c) $T = \ln(\text{Mg/Ca}/0.4)/0.09$:
$\ln 5/0.09 = \qty{17.9}{\celsius}$ dan $\ln 2.5/0.09 = \qty{10.2}{\celsius}$.
\end{solution}

\begin{solution}{exo:b2:climate-proxies:4}
(a) $2668 - 2213 + 1 = 456$ tahun. (b) Satu pohon dapat kehilangan satu lingkaran secara lokal, atau menunjukkan lingkaran
yang tidak dapat dibaca; mencocokkan pola lingkaran lebar dan sempit pada banyak pohon mengungkap
lingkaran yang hilang dan menetapkan tahun setiap lingkaran, dan kesesuaian antarpohon mengukur
keandalan penanggalannya. (c) Di batas pohon satu faktor, yaitu dinginnya musim tumbuh,
membatasi semua pohon, sehingga pohon-pohon itu berbagi pola tahun baik dan tahun buruk yang sama; di dataran rendah yang beriklim lunak dan
lembap tidak ada satu faktor pun yang dominan, pola pohon-pohonnya berbeda, dan \omterm{def:b2:climate-proxies:dendro}{penanggalan silang} dapat gagal.
\end{solution}

\begin{solution}{exo:b2:climate-proxies:5}
(a) $125 \times 40/(3682.2 - 125) = 5000/3557.2 = 1.41$ permil; $134 \times 40/3548.2 =
5360/3548.2 = 1.51$ permil. (b) Dengan $H - h = 3552.2$: $3900/3552.2 = 1.10$, $4550/3552.2 = 1.28$,
dan $6500/3552.2 = 1.83$ permil. (c) $\delta_{\mathrm{ice}} = -1.0 \times 3552.2/130 = -27.3$
permil, lebih berat daripada lembaran es masa kini mana pun (Greenland $-30$ sampai $-35$, lembaran es Amerika Utara
$-40$, es Antarktika serendah $-50$). Neraca massa dengan es yang realistis dan pengukuran air
pori tidak sesuai; selisih sekitar setengah permil itu adalah ketidakpastian nyata metode ini.
\end{solution}

\begin{solution}{exo:b2:climate-proxies:6}
(a) $0.372\,\mathrm{e}^{0.0936T} = 0.4\,\mathrm{e}^{0.09T}$, jadi $\mathrm{e}^{0.0036T} =
0.4/0.372 = 1.0753$, sehingga $T = \ln 1.0753/0.0036 = \qty{20.2}{\celsius}$. (b) Pada
\qty{2}{\celsius}: $0.4\,\mathrm{e}^{0.18} = 0.479$ dan $0.372\,\mathrm{e}^{0.1872} = 0.449$, sehingga
rasio hasil pencocokan adalah $0.449/0.479 = 0.937$, \qty{6.3}{\percent} lebih rendah. (c) $T = \ln(0.449/0.4)/0.09 =
\ln 1.1225/0.09 = \qty{1.3}{\celsius}$: hukum ramalan membaca cangkang itu sekitar
\qty{0.7}{\celsius} terlalu dingin. Di air yang dingin, pilihan kalibrasi berpengaruh sampai tingkat satu
derajat.
\end{solution}

\begin{solution}{exo:b2:climate-proxies:7}
(a) Dari Holosen sampai GI-1, \qty{112.19}{m} es untuk 2989 tahun: $11\,219/2989 =
\qty{3.75}{cm}$ per tahun; dari GI-16 sampai GI-17, \qty{17.89}{m} untuk 1160 tahun: $1789/1160 =
\qty{1.54}{cm}$ per tahun. (b) Es
mengalir: ketika setiap lapisan tertimbun, lapisan itu menyebar ke samping dan menipis, dan lapisan yang lebih dalam dan lebih tua
telah lebih lama menipis. (c) Bila lapisan menjadi terlalu tipis, satu sentimeter atau kurang, instrumen
tidak lagi dapat membedakan satu tahun dari tahun berikutnya, dan \omterm{def:b2:climate-proxies:age-model}{model umur} harus memakai titik ikat atau model aliran
sebagai gantinya.
\end{solution}

\begin{solution}{exo:b2:climate-proxies:8}
(a) Bagian selang itu adalah $(1550 - 1492.45)/(1604.64 - 1492.45) = 57.55/112.19 = 0.513$,
sehingga umurnya $11\,703 + 0.513 \times 2989 = 11\,703 + 1533 = \num{13236}$ tahun b2k. (b) Tahun 1950 adalah 50
tahun sebelum 2000: \num{13186} tahun sebelum 1950. (c) $50 + 0.513 \times (93 - 50) = 72$ tahun.
\end{solution}

\begin{solution}{exo:b2:climate-proxies:9}
(a) Setiap lapisan tidak pasti menyumbang setengah tahun: $2 \times 93 = 186$ dan $2 \times 1287 = 2574$
lapisan tidak pasti. (b) $\frac12\sqrt{186} = 6.8$ tahun dan $\frac12\sqrt{2574} = 25.4$ tahun.
(c) $93/6.8 = 14$ dan $1287/25.4 = 51$. Perbandingan $\frac12 n$ terhadap $\frac12\sqrt{n}$ adalah
$\sqrt{n}$: perbandingan itu tumbuh sebagai akar kuadrat jumlah lapisan tidak pasti, sehingga pilihan antara
kedua aturan itu makin penting seiring kedalaman.
\end{solution}

\begin{solution}{exo:b2:climate-proxies:10}
(a) Pada \omterm{prop:b2:isotope-geochemistry:mwl}{garis air meteorik} $\delta$D berubah delapan kali sebanyak $\delta^{18}$O, sehingga
kemiringannya $8 \times 0.8 = 6.4$ permil $\delta$D per derajat. Penurunannya $441.7 - 396.6 = 45.1$
permil, jadi pendinginannya $45.1/6.4 = \qty{7.0}{\celsius}$. (b) Samudra glasial sekitar
$8 \times 1.0 = 8$ permil lebih berat dalam $\delta$D; salju yang terbentuk dari samudra masa kini akan
lebih ringan sebanyak itu, $-441.7 - 8 = -449.7$ permil. Penurunannya menjadi $53.1$ permil dan pendinginannya
$53.1/6.4 = \qty{8.3}{\celsius}$. (c) $9.39 - 0.09 = \qty{9.30}{\celsius}$, sehingga $45.1/9.30 = 4.85$
permil per derajat. Deret yang diterbitkan juga dikoreksi untuk ketinggian lembaran es yang
berubah, dan hubungan suhunya berasal dari model yang menyimulasikan isotop air, bukan dari
kemiringan spasial saja.
\end{solution}

\begin{solution}{exo:b2:climate-proxies:11}
(a) Logaritma rasio-rasio itu adalah $-0.528$, $-0.223$, $0.191$, $0.501$, $0.904$, $1.227$;
rata-ratanya $0.345$ dan suhu rata-ratanya 14. Dengan $S_{TT} = \sum (T - 14)^2 = 280$ dan $S_{T\ell} =
\sum (T - 14)(\ell - 0.345) = 24.93$, $A = 24.93/280 = \qty{0.0890}{\per\celsius}$ dan $\ln B =
0.345 - 0.0890 \times 14 = -0.901$, sehingga $B = \qty{0.406}{mmol/mol}$. (b) Residunya adalah
$+0.017$, $-0.034$, $+0.023$, $-0.023$, $+0.025$, dan $-0.009$, beberapa persen dari rasio itu, tanpa
tren; hasil pencocokan dekat dengan $0.4\,\mathrm{e}^{0.09T}$. (c) $T = (\ln 1.30 + 0.901)/0.0890 = (0.262 +
0.901)/0.0890 = \qty{13.1}{\celsius}$. (d) $T = (\ln 0.45 + 0.901)/0.0890 = (-0.799 +
0.901)/0.0890 = \qty{1.2}{\celsius}$. Nilai ini terletak di bawah sampel kalibrasi yang paling dingin,
\qty{4}{\celsius}: hukum itu diekstrapolasi, dan pada suhu rendah Mg/Ca cangkang diketahui
menyimpang dari hukum eksponensial di bawah pengaruh ion karbonat.
\end{solution}

\begin{solution}{exo:b2:climate-proxies:12}
(a) Per satuan luas samudra, air samudra masa kini, $H$ pada 0, terbagi antara simpanan-simpanan
$h_i$ pada $\delta_i$ dan samudra yang tersisa $H - h$, dengan $h = \sum h_i$: $0 = \sum h_i\delta_i +
(H - h)\,\Delta\delta$. Setelah tahap pertama $\Delta\delta = 60 \times 30/(3682.2 - 60) =
1800/3622.2 = 0.50$ permil; setelah tahap kedua $(1800 + 70 \times 45)/(3682.2 - 130) =
4950/3552.2 = 1.39$ permil. (b) $-4950/130 = -38.1$ permil. (c) Aturan itu adalah $1.39/130 = 0.0107$
permil per meter, dan memberikan $0.50/0.0107 = \qty{47}{m}$ untuk tahap pertama alih-alih
\qty{60}{m}: es pertama lebih berat daripada rata-rata, sehingga setiap meternya menggeser samudra
lebih sedikit. Aturan linear yang dikalibrasi pada satu keadaan salah membaca keadaan lain bila komposisi rata-rata
es berubah.
\end{solution}

\begin{solution}{pb:b2:climate-proxies:1}
\textbf{1.} $361.84 \times 10^6\ \unit{km^2} \times \qty{0.130}{km} = \qty{4.70e7}{km^3}$, sekitar 47
juta kilometer kubik.
\textbf{2.} $130 \times 40/(3682.2 - 130) = 5200/3552.2 = 1.46$ permil.
\textbf{3.} $5000/3557.2 = 1.41$ dan $5360/3548.2 = 1.51$ permil.
\textbf{4.} $4550/3552.2 = 1.28$ dan $6500/3552.2 = 1.83$ permil. Komposisi es
menggeser jawaban sebesar \num{0.55} permil, muka laut sebesar \num{0.10}: es lebih penting.
\textbf{5.} $\delta_{\mathrm{ice}} = -1.0 \times 3552.2/130 = -27.3$ permil.
\textbf{6.} $40/3682.2 = 0.0109$ permil per meter, yang dibulatkan menjadi \num{0.011} dari muka laut
yang bertanggal.
\textbf{7.} Pada $\delta^{18}\mathrm{O}_{\mathrm{sw}}$ tetap, $\mathrm{d}T = -4.76\,\mathrm{d}\delta^{18}
\mathrm{O}_{\mathrm{b}}$, sehingga $\mathrm{d}T = -1$ memberikan $\mathrm{d}\delta^{18}\mathrm{O}_{\mathrm{b}} =
1/4.76 = 0.21$ permil.
\textbf{8.} $\delta^{18}\mathrm{O}_{\mathrm{b}} = (16.1 - T)/4.76 + \delta^{18}\mathrm{O}_{\mathrm{sw}}
- 0.27 = 14.1/4.76 - 0.27 = 2.96 - 0.27 = 2.69$ permil.
\textbf{9.} $17.1/4.76 + 1.0 - 0.27 = 3.59 + 0.73 = 4.32$ permil. Perubahannya $4.32 - 2.69 =
1.63$ permil: \num{1.00} dari air dan $3.0 \times 0.21 = 0.63$ dari pendinginan sebesar
\qty{3}{\celsius}.
\textbf{10.} $T = \ln(2.0/0.4)/0.09 = \ln 5/0.09 = \qty{17.9}{\celsius}$.
\textbf{11.} Pendinginan sebesar $0.18/0.1 = \qty{1.8}{\celsius}$.
\textbf{12.} Mg/Ca adalah perbandingan dua unsur dalam kalsit, yang ditentukan oleh suhu
kalsifikasi; perbandingan itu tidak melibatkan isotop oksigen, sehingga tidak terpengaruh oleh komposisi
air yang diubah oleh es.
\textbf{13.} $4.92 - 3.23 = 1.69$ permil.
\textbf{14.} Dengan menyelesaikan $\Delta = h \times 40/(H - h)$ untuk $h$: $h = H\Delta/(\Delta + 40) = 3682.2
\times 1.69/41.69 = \qty{149}{m}$, lebih dari 125 sampai \qty{134}{m} menurut bukti muka laut:
sebagian perubahan itu pasti berupa pendinginan.
\textbf{15.} $1.69 - 1.46 = 0.23$ permil, dan $4.76 \times 0.23 = \qty{1.1}{\celsius}$.
\textbf{16.} $1.69 - 1.00 = 0.69$ permil, dan $4.76 \times 0.69 = \qty{3.3}{\celsius}$ (3.28).
\textbf{17.} $0.23/0.25 = \qty{0.9}{\celsius}$ dan $0.69/0.25 = \qty{2.8}{\celsius}$.
\textbf{18.} Air Dalam Atlantik Utara mendingin dari 2--\qty{4}{\celsius} ke \qty{-1.1}{\celsius},
sebesar 3.1 sampai \qty{5.1}{\celsius}; Air Dasar Antarktika dari 0--\qty{1}{\celsius} ke
\qty{-1.9}{\celsius}, sebesar 1.9 sampai \qty{2.9}{\celsius}; Mg/Ca Atlantik Barat Laut memberikan
\qty{1.8}{\celsius}. Semuanya terletak antara sekitar 2 dan \qty{5}{\celsius}, di sekitar 2.8 sampai
\qty{3.3}{\celsius} dari koreksi air pori, dan di atas \qty{1}{\celsius} dari neraca
massa dengan es pada $-40$ permil.
\textbf{19.} $-396.6 - (-441.7) = 45.1$ permil.
\textbf{20.} $8 \times 0.8 = 6.4$ permil per derajat, sehingga $45.1/6.4 = \qty{7.0}{\celsius}$.
\textbf{21.} $9.39 - 0.09 = \qty{9.30}{\celsius}$, dan $45.1/9.30 = 4.85$ permil per derajat.
\textbf{22.} Dome C, \qty{9.3}{\celsius}, lebih dari permukaan global, 5 sampai \qty{7}{\celsius}, yang lebih
dari samudra dalam, sekitar \qty{3}{\celsius}.
\textbf{23.} Air dalam sudah dingin pada masa kini, 0 sampai \qty{4}{\celsius}, dan air laut tidak dapat
mendingin jauh di bawah titik bekunya, yang hampir dicapai air dasar glasial pada
\qty{-1.9}{\celsius}: samudra dalam hanya dapat kehilangan beberapa derajat, sedangkan udara di atas Antarktika
tidak memiliki batas bawah semacam itu.
\textbf{24.} Dengan koreksi air pori, samudra dalam sekitar \textbf{\qty{3.3}{\celsius}
lebih dingin} pada Maksimum Glasial Terakhir dibandingkan dengan masa kini, dekat dengan 2 sampai \qty{5}{\celsius} dari
perkiraan-perkiraan air dalam yang independen.
\end{solution}
